% Удиртгал

\phantomsection
\addchaptertocentry{Удиртгал}
\section*{Удиртгал}
Бүтээмжийг нэмэгдүүлэх нь аж ахуйн нэгж, төрийн байгууллагын өрсөлдөх чадварын гол хүчин зүйл юм. Олон байгууллага 5S, Kaizen, lean зэрэг аргачлалыг нэвтрүүлж байгаа ч эдгээр аргачлал нэвтэрсэн байгууллагуудад шалгах хуудас цаасаар, ажлын даалгавар ярианы хэлбэрээр, тайлан Excel хүснэгтээр тус тусдаа хөтлөгддөг. Үүний улмаас аудитаар илэрсэн дутагдал засагдсан эсэхийг хэн ч хянадаггүй, менежер ажлын явцыг бодит хугацаанд хардаггүй, сарын тайланг олон эх сурвалжаас гараар нэгтгэхэд олон цаг зарцуулдаг байна.

Иймээс 5S болон бүтээмжийн хэрэгслүүдийг өдөр тутмын ажлын удирдлагатай нэг системд холбож, тайланг автоматаар бүрдүүлдэг, монгол хэл дээр ажилладаг, интернэт тасалдалтай цех, агуулахад ч ашиглах боломжтой платформ хэрэгцээтэй байна.

\section*{Зорилго}
Энэхүү дипломын ажлын зорилго нь байгууллагын 5S, бүтээмжийн хэрэгслүүд болон ажлын удирдлагыг нэгтгэж, менежерт ажлыг бодит хугацаанд оноож хянах, ажилтанд ажлаа гар утаснаас хөтлөх, тайланг автоматаар гаргах боломж олгох вэб болон гар утасны платформыг зохион бүтээж, хэрэгжүүлэхэд оршино.

\section*{Зорилт}
Зорилгод хүрэхийн тулд дараах зорилтуудыг тавьсан.
\begin{itemize}
    \item 5S, Gemba, Kaizen, шатлалт өдөр тутмын хурал зэрэг аргачлалын онолыг судлах;
    \item ижил төстэй системүүдийг харьцуулан шинжилж, давуу болон сул талыг тодорхойлох;
    \item 5S, lean аргачлал хэрэглэдэг байгууллагуудын ажилтнуудаас хэрэглэгчийн судалгаа авах;
    \item функциональ болон функциональ бус шаардлагыг тодорхойлох;
    \item UML диаграммууд, системийн архитектур, өгөгдлийн сангийн схем, гол алгоритмуудыг зохиомжлох;
    \item backend API, вэб апп, гар утасны аппыг хөгжүүлж, автомат тестээр баталгаажуулах;
    \item туршилтын байгууллагад ажиллуулж, хэрэглэхэд хялбар байдлыг SUS аргаар үнэлэх;
    \item төслийн эдийн засгийн үр ашгийг тооцох.
\end{itemize}

\section*{Систем хөгжүүлэх үндэслэл}
\begin{enumerate}[nosep]
    \item 5S аудитаар илэрсэн дутагдлыг хариуцагчтай, хугацаатай ажил болгох механизм байхгүй тул аудит нь хэмжилт төдий үлддэг.
    \item Менежер ажлын явц, хоцролтыг бодит хугацаанд харах боломжгүй, хоцорсон ажил хурал дээр л илэрдэг.
    \item Сарын, хагас жилийн, жилийн тайланг гараар нэгтгэдэг нь цаг их зарцуулж, алдаа гаргах эрсдэлтэй.
    \item Ажилчдын сайжруулалтын санаа бүртгэгддэггүй, санал гаргасан хүн хариугаа авдаггүй.
    \item Гадаадын ижил төстэй системүүд монгол хэлгүй, хэрэглэгч тус бүрээр гадаад валютаар төлбөртэй, байгууллагын өөрийн серверт байршуулах боломжгүй.
\end{enumerate}

\section*{Судлагдсан байдал}
Lean удирдлага, 5S аргачлалыг Хирано~\cite{hirano}, Имай~\cite{imai}, Лайкер~\cite{liker} нар системтэйгээр боловсруулсан бөгөөд тэдгээрийг цахимжуулсан SafetyCulture, KaiNexus, Tervene, Weekdone зэрэг гадаадын системүүд байдаг. Эдгээрийг нэгдүгээр бүлгийн 1.2-т дэлгэрэнгүй харьцуулав. Дотоодод 5S болон бүтээмжийн хэрэгслүүдийг ажлын удирдлагатай нэгтгэсэн, монгол хэл дээрх систем олдсонгүй.

\section*{Шинэлэг тал}
\begin{itemize}
    \item 5S аудитын оноо босгоос доош бол засах ажил талбайн хариуцагчид автоматаар оноогдож, мэдэгдэл очно;
    \item өглөөний хурал, Gemba явалт, долоо хоногийн тайлан зэрэг lean удирдлагын дадлыг нэг системд холбосон;
    \item сарын тайлан бүртгэлээс автоматаар бүрдэж, тогтсон өдөр хаагдан өөрчлөгдөхгүйгээр архивлагдана;
    \item гар утасны апп сүлжээгүй үед хийсэн бүртгэлийг хадгалж, сүлжээ сэргэхэд давхардуулалгүйгээр илгээнэ;
    \item бүрэн монгол, англи хэлтэй, WCAG 2.1 AA хүртээмжийн шаардлагыг хангахаар шалгагдсан.
\end{itemize}

\section*{Ач холбогдол}
Систем нь аудит, ажил, тайлангийн мэдээллийг үүссэн газарт нь бүртгэж, нэг өгөгдлийн санд нэгтгэснээр менежерийн шийдвэр гаргалтыг бодит мэдээлэлд тулгуурлуулна. Тайлан бэлтгэх, ажлын явц тодруулах цагийг хэмнэж, 5S-ийг нэг удаагийн кампанит ажлаас тогтвортой дадал болгоход дэмжлэг үзүүлнэ.

\section*{Хамрах хүрээ}
Систем нь тодорхой нэг байгууллагад бус, 5S, lean аргачлал хэрэглэдэг аль ч байгууллагад зориулагдсан. Туршилтын ажиллагааг Монголын Бүтээмжийн төв (МБТ)-д хийхээр төлөвлөсөн. Систем олон байгууллагад зэрэг үйлчлэх (multi-tenant) бүтэцтэй тул 5S, lean аргачлал нэвтрүүлж буй үйлдвэр, үйлчилгээ, төрийн байгууллагуудад өргөтгөх боломжтой. Санхүүгийн бүртгэл, цалин, хүний нөөцийн бүрэн систем хамрах хүрээнд орохгүй.
